%% evaluation.tex  —  Moufida Evaluation Framework and Results
%% pdflatex evaluation.tex

\documentclass[11pt,a4paper]{article}

\usepackage[margin=2.2cm]{geometry}
\usepackage{amsmath,amssymb}
\usepackage{graphicx}
\usepackage[hidelinks]{hyperref}
\usepackage{xcolor}
\usepackage{booktabs}
\usepackage{enumitem}
\usepackage{tcolorbox}
\usepackage{microtype}
\usepackage{titlesec}
\tcbuselibrary{skins,breakable}

%% ── Palette (section titles and box borders/backgrounds only) ───────────────
\definecolor{EvalGreen}{RGB}{27,115,65}
\definecolor{EvalTeal}{RGB}{0,100,88}
\definecolor{EvalOrange}{RGB}{172,85,0}
\definecolor{EvalSlate}{RGB}{50,62,75}
\definecolor{PassBg}{RGB}{232,248,238}
\definecolor{PendingBg}{RGB}{255,244,224}
\definecolor{MetricBg}{RGB}{236,245,245}
\definecolor{AdminBg}{RGB}{245,246,248}
\definecolor{BoxBg}{RGB}{250,252,250}

%% ── Status badges (black text on coloured background) ───────────────────────
\newcommand{\ppass}{\colorbox{PassBg}{\small\textbf{PASS}}}
\newcommand{\ppending}{\colorbox{PendingBg}{\small\textbf{PENDING}}}

%% ── Spacing ─────────────────────────────────────────────────────────────────
\setlength{\parskip}{4pt plus 1pt}
\setlength{\parindent}{0pt}
\setlist{itemsep=2pt, topsep=3pt, parsep=0pt, leftmargin=1.4em}
\setlist[enumerate]{itemsep=2pt, topsep=3pt}

%% ── Section styling ─────────────────────────────────────────────────────────
\titlespacing*{\section}{0pt}{14pt}{5pt}
\titlespacing*{\subsection}{0pt}{10pt}{3pt}
\titleformat{\section}{\large\bfseries\color{EvalGreen}}{{\thesection}}{0.6em}{}[\vspace{-2pt}\color{EvalGreen}\titlerule]
\titleformat{\subsection}{\normalsize\bfseries\color{EvalSlate}}{{\thesubsection}}{0.6em}{}

%% ── Colour boxes ────────────────────────────────────────────────────────────
\tcbset{
  base/.style={
    boxrule=0.6pt, arc=3pt,
    boxsep=4pt, top=4pt, bottom=4pt, left=6pt, right=6pt,
    coltitle=white, fonttitle=\small\bfseries, breakable
  }
}
\newtcolorbox{passbox}[1]{base,
  colback=PassBg, colframe=EvalGreen!75,
  title={#1}}
\newtcolorbox{pendingbox}[1]{base,
  colback=PendingBg, colframe=EvalOrange!75,
  title={#1}}
\newtcolorbox{metricbox}[1]{base,
  colback=MetricBg, colframe=EvalTeal!70,
  title={#1}}
\newtcolorbox{adminbox}[1]{base,
  colback=AdminBg, colframe=EvalSlate!50,
  title={#1}}

%% ── Title ───────────────────────────────────────────────────────────────────
\title{%
  {\color{EvalGreen}\rule{\linewidth}{1.5pt}}\\[6pt]
  \textbf{\Large Evaluation Framework and Results}\\[4pt]
  {\normalsize Moufida Startup Diagnostic Platform}\\[6pt]
  {\color{EvalGreen}\rule{\linewidth}{0.8pt}}
}
\author{Team Makrouna Kadheba}
\date{AINS 2026}

\begin{document}
\maketitle\thispagestyle{empty}

\begin{abstract}\small
Evaluating a multi-component AI diagnostic system requires more than a held-out test set.
This document describes Moufida's three-tier evaluation framework:
\textbf{T2} (component-level invariants — determinism, stability, anomaly recall),
\textbf{T1} (end-to-end semantic quality — maturity stage classification), and
\textbf{T3} (retrieval quality — RAG recall and ranking).
Subsequent sections cover the CBM calibration loop as a continuous self-evaluation
mechanism, the Score Debate feature as a human-in-the-loop correction pathway, and
the admin observability panel as a live evaluation surface during demonstrations.
\end{abstract}

\tableofcontents
\vspace{6pt}
\noindent\color{EvalGreen}\rule{\linewidth}{0.4pt}

%% ─────────────────────────────────────────────────────────────────────────────
\section{Evaluation Philosophy}
%% ─────────────────────────────────────────────────────────────────────────────

Moufida is not a classification model to be benchmarked against a static dataset.
It is a pipeline of heterogeneous components — a deterministic scoring engine,
an LLM-powered rubric system, a hybrid retrieval pipeline, a CAT intake engine,
and a daemon that observes real-world signals. Each component requires a different
evaluation lens.

\begin{metricbox}{Why standard ML evaluation is insufficient here}
\begin{itemize}[leftmargin=1em]
  \item The scoring engine must be \textbf{deterministic}: a founder who runs a diagnosis
        twice must get the same result.
  \item The LLM rubric must be \textbf{stable}: small temperature fluctuations must not
        produce actionably different scores.
  \item The anomaly detectors must have \textbf{high recall}: a contradictory profile that
        passes silently is worse than a false positive.
  \item The RAG pipeline must retrieve the \textbf{right evidence}: a Tunisian founder
        asking about grants should get APII and BFPME, not generic EU startup advice.
  \item The CAT intake must \textbf{converge efficiently}: reaching a low standard error
        in fewer questions than a fixed-item survey.
\end{itemize}
\end{metricbox}

\subsection{Three-Tier Structure}

\begin{center}\small
\begin{tabular}{@{}llll@{}}
\toprule
\textbf{Tier} & \textbf{Scope} & \textbf{Primary metric} & \textbf{Status} \\
\midrule
T2 & Component properties (determinism, stability, anomaly) & Pass rate, $\sigma$ & \ppass \\
T1 & End-to-end semantic quality (maturity classification)  & Macro-F1, Top-2     & \ppending \\
T3 & Retrieval quality (RAG recall and ranking)             & Recall@3, MRR       & \ppending \\
\bottomrule
\end{tabular}
\end{center}

T2 tests are runnable without any labelled dataset — they verify invariants that must
hold regardless of the startup profiles used.
T1 and T3 require ground-truth annotations currently in preparation.

%% ─────────────────────────────────────────────────────────────────────────────
\section{T2 — Component-Level Evaluation}
%% ─────────────────────────────────────────────────────────────────────────────

T2 tests live in \texttt{backend/eval/} and can be run against any live stack.
They are property-based: they assert invariants rather than compare against ground truth.

\subsection{T2a — Scoring Determinism \ppass}

\begin{passbox}{T2a Result: 10\,/\,10 runs identical across all axes and sub-dimensions}
\textbf{What we tested.}
The Affinitree scoring engine must produce bit-for-bit identical results given the same
profile and evidence tiers, across all 10 axis scores, all sub-dimension breakdowns,
and all CBM concept scores.

\textbf{Protocol.}
\begin{enumerate}
  \item Fix one synthetic profile with all fields at the T2 evidence tier.
  \item Run the full scoring pipeline (axes + CBM pass) 10 consecutive times.
  \item Compare all 10 output objects pairwise for exact equality.
\end{enumerate}
\end{passbox}

\begin{itemize}
  \item The scoring formula $S = \sum w_i \cdot v_i \cdot m_i$ is a closed-form Rust
        computation. The same inputs always follow the same code path, with no floating-point
        non-determinism.
  \item The LLM concept-scoring call uses \texttt{temperature=0} and \texttt{seed=42}.
        Both are required simultaneously: temperature 0 alone does not guarantee
        determinism across all Ollama backends.
  \item This guarantee is load-bearing for the \textbf{Score Debate} feature. A human
        expert who locks an override can always compare it against the deterministic
        model baseline — if the baseline shifted randomly, the comparison would be meaningless.
  \item It also makes regression testing practical: a code change that accidentally shifts
        a score is immediately detectable by re-running T2a.
\end{itemize}

\subsection{T2b — Rubric Stability \ppending}

\begin{pendingbox}{T2b Target: $\sigma \leq 0.15$ per field on [0,\,4] scale
— preliminary $\bar{\sigma} = 0.07$}
\textbf{What we tested.}
When the LLM rubric scoring runs under realistic temperature (\texttt{temperature=0.3}),
the per-field variance across 10 runs must be small enough that no actionable advice changes.

\textbf{Protocol.}
\begin{enumerate}
  \item Select 3 profiles at different maturity stages (ideation, MVP, growth).
  \item Run rubric scoring 10 times per profile with \texttt{temperature=0.3}.
  \item Compute $\sigma$ per field; report $\sigma_\text{max}$ and $\bar{\sigma}$.
\end{enumerate}

\textbf{Preliminary result} (3 profiles, 3 fields):
$\sigma_\text{max} = 0.11$, $\bar{\sigma} = 0.07$.
Full run pending (5 profiles, all rubric fields).
\end{pendingbox}

\begin{itemize}
  \item The threshold $\sigma \leq 0.15$ on a [0,\,4] scale translates to $<2$ points of
        variance on the final [0,\,100] axis score. A score that moves by $<2$ points
        between runs does not change the bottleneck identification or roadmap recommendations.
  \item Temperature 0.3 was chosen as the realistic production setting: a small amount of
        variability is desirable so that different founders with similar profiles do not
        receive word-for-word identical justifications.
  \item Preliminary results ($\bar{\sigma} = 0.07$) suggest the system is well within
        target. The remaining question is whether outlier fields exist — hence the full run.
\end{itemize}

\subsection{T2c — Anomaly Detection Recall \ppass}

\begin{passbox}{T2c Result: 10\,/\,10 planted contradictions detected}
\textbf{What we tested.}
The 10 rule-based anomaly detectors must collectively flag every major logical
contradiction that can appear in a startup profile.

\textbf{Protocol.}
\begin{enumerate}
  \item Construct 10 synthetic profiles, each with exactly one planted contradiction:
  \begin{itemize}
    \item Revenue claimed with no product or MVP
    \item Large team claimed but no hiring evidence
    \item High market score with no market research cited
    \item IP protection claimed with no legal filings
    \item Growth-stage self-assessment with ideation-stage evidence
    \item Sales figures cited with zero customer interviews
    \item International target market with no export compliance steps
    \item Grant applications cited but no registered legal entity
    \item Burn rate inconsistent with stated runway
    \item Competitor analysis present but no competitor named
  \end{itemize}
  \item Run each profile through the full diagnostic pipeline.
  \item Verify that the planted flag appears in the \texttt{anomalies} field of
        \texttt{diagnostic\_history}.
\end{enumerate}
\end{passbox}

\begin{itemize}
  \item Anomaly detection is itself an interpretability mechanism, not just a quality
        check: a profile that scores 72/100 on market-intelligence while flagging
        ``no market research cited'' is numerically plausible but logically suspect.
        The flag surfaces this contradiction to the founder and any investor reading the report.
  \item Rule-based detectors were chosen over statistical anomaly detection because they
        are \emph{auditable}: each rule can be read by a non-technical reviewer, and each
        flag message explains exactly which contradiction was found.
  \item The detectors operate at the orchestrator level, after all axis scores are collected,
        so they have full profile visibility to compare cross-axis claims.
\end{itemize}

%% ─────────────────────────────────────────────────────────────────────────────
\section{T1 — End-to-End Semantic Evaluation}
%% ─────────────────────────────────────────────────────────────────────────────

\begin{pendingbox}{T1 Target: Macro-F1 $\geq 0.65$, Top-2 accuracy $\geq 0.85$
— preliminary 0.71\,/\,0.92}
\textbf{Task.}
Given a synthetic startup profile, the orchestrator must classify it into one of five
maturity stages: \{ideation, pre-seed, seed, growth, scale\}.

\textbf{Ground truth.}
50 synthetic profiles with stage labels assigned by a domain expert (Tunisian startup
ecosystem practitioner).
Dataset construction in progress (25/50 profiles complete).

\textbf{Protocol.}
\begin{enumerate}
  \item Run all 50 profiles through the full diagnostic pipeline.
  \item Compare the predicted stage to the expert label.
  \item Compute per-class F1 and macro-average (equally weighted across 5 classes).
  \item Compute Top-2 accuracy: the prediction is correct if within one stage of the label.
\end{enumerate}

\textbf{Early result} (25 profiles): Macro-F1 = 0.71, Top-2 accuracy = 0.92.
\end{pendingbox}

\begin{itemize}
  \item Macro-F1 (not micro) is the primary metric because the stage distribution is
        naturally imbalanced: there are far more ideation and pre-seed profiles than
        growth or scale. Macro-F1 gives equal weight to each class, preventing the
        metric from being dominated by the majority class.
  \item Top-2 accuracy is the more operationally relevant metric: a prediction one stage
        off produces a roadmap that is still coherent and only marginally mis-calibrated.
        A two-stage error would produce meaningfully wrong advice.
  \item \textbf{Perception Gap} is an auxiliary metric derived from T1. On the 25-profile
        sample, 14/25 profiles exhibited a mismatch between the computed stage and the
        founder's self-assessment — consistent with domain expertise that early-stage
        founders systematically over-estimate their maturity.
\end{itemize}

%% ─────────────────────────────────────────────────────────────────────────────
\section{T3 — Retrieval Quality}
%% ─────────────────────────────────────────────────────────────────────────────

\begin{pendingbox}{T3 Target: Recall@3 $\geq 0.80$, MRR $\geq 0.70$
— preliminary 0.90\,/\,0.83}
\textbf{Task.}
Given a (query, relevant document set) pair, verify that the hybrid RAG pipeline
retrieves the correct documents in the top-3 results.

\textbf{Query set.}
30 manually constructed queries covering typical axis retrieval scenarios —
e.g., ``StartupAct fiscal advantages'', ``BFPME loan requirements'',
``IP registration steps in Tunisia'', ``minimum capital for SA in Tunisia''.

\textbf{Metrics.}
\begin{itemize}
  \item \textbf{Recall@3}: fraction of relevant documents appearing in the top-3 results.
  \item \textbf{MRR (Mean Reciprocal Rank)}:
        $\dfrac{1}{|Q|}\sum_{q} \dfrac{1}{\text{rank of first relevant result for } q}$
\end{itemize}

\textbf{Early result} (10 queries, KB only, no web search): Recall@3 = 0.90, MRR = 0.83.
\end{pendingbox}

\begin{itemize}
  \item Recall@3 was chosen over Recall@1 because the pipeline feeds 5 chunks to each
        axis LLM call. Requiring the relevant document in the top 3 is a stricter-than-necessary
        bar and gives a conservative lower bound on actual evidence quality.
  \item \textbf{Axis direction probe contribution.} On the 10-query preliminary run,
        disabling the probe re-ranking (reverting to plain BM25+dense RRF) dropped
        Recall@3 from 0.90 to 0.70 and MRR from 0.83 to 0.65.
        This $+20$ pp Recall improvement directly validates the probe's contribution for
        axis-specific queries.
  \item The query set deliberately includes cross-lingual queries (French and Darija input,
        French and Arabic documents) to test bge-m3's multilingual alignment.
\end{itemize}

%% ─────────────────────────────────────────────────────────────────────────────
\section{Continuous Self-Evaluation via CBM Calibration}
%% ─────────────────────────────────────────────────────────────────────────────

\subsection{The Calibration Loop as a Feedback Mechanism}

Every time a human reviewer overrides a score and locks it in the Score Debate UI,
the locked value becomes a training signal for the CBM weight calibration:

\begin{metricbox}{Ridge Calibration — Continuous Weight Estimation}
\[
  \hat{\mathbf{w}}^{(t)} \;=\; \arg\min_{\mathbf{w}}\;
  \bigl\| X^{(t)}\mathbf{w} - \mathbf{y}^{(t)} \bigr\|_2^2
  \;+\; \lambda\,\|\mathbf{w}\|_2^2
\]
$X^{(t)} \in \mathbb{R}^{n_t \times k}$ grows with each diagnostic run;
$\mathbf{y}^{(t)}$ grows with each locked override.
\end{metricbox}

\begin{itemize}
  \item This makes the CBM a \textbf{self-improving} evaluation mechanism: the more
        founders use the system and the more overrides are locked, the more concept
        weights reflect ground truth in the Tunisian startup context.
  \item Calibration quality is tracked via two internal metrics: weight stability
        $\|\hat{\mathbf{w}}^{(t)} - \hat{\mathbf{w}}^{(t-1)}\|_2$ (should decrease
        toward zero as the system converges) and residual error on held-out locked
        snapshots.
  \item After 20+ locked overrides, preliminary runs show
        $\|\hat{\mathbf{w}}^{(t)} - \hat{\mathbf{w}}^{(t-1)}\| < 0.01$,
        indicating weights have stabilised around a consistent interpretation of
        concept importance for Tunisian startups.
\end{itemize}

\subsection{Score Debate and Regression Prevention}

\begin{itemize}
  \item The Score Debate feature lets a human reviewer challenge any axis score.
        The panel presents the full CBM concept breakdown, the retrieved evidence chunks
        with axis-relevance scores, and the LLM justification.
  \item If the reviewer overrides, the new score is saved with \texttt{locked = true}
        in \texttt{score\_snapshots}. A locked score cannot be downgraded by a subsequent
        automatic diagnostic — the engine detects the lock and skips re-scoring.
  \item This is critical when Moufida reports are shared with investors or grant bodies:
        an expert-validated score must not silently revert to a lower value on the next run.
\end{itemize}

\subsection{Bottleneck Stability}

The same profile is run 10 times with \texttt{temperature=0.3} (allowing LLM variability)
and the bottleneck concept is recorded each time.

Result: \textbf{9/10 runs identify the same bottleneck concept.}

This means that even under realistic temperature conditions, the explanation is stable —
the founder reliably sees the same actionable recommendation across runs.

%% ─────────────────────────────────────────────────────────────────────────────
\section{Admin Observability Panel as Live Evaluation}
%% ─────────────────────────────────────────────────────────────────────────────

The admin panel (\texttt{admin/}, port 3002) provides a real-time window into every
evaluation-relevant event during a demonstration or audit.

\begin{adminbox}{Distributed Request Tracing}
Every HTTP request receives a \texttt{request\_id} (UUID) propagated through all downstream
calls via a Python ContextVar and stored in \texttt{api\_requests} and \texttt{llm\_calls}.
Clicking Trace on any request row reveals the full call tree:
\begin{itemize}[leftmargin=1em]
  \item Which axis services were called and in which wave
  \item Which Ollama calls followed: model, prompt hash, duration, token counts
  \item Which RAG retrievals occurred: query text, number of results returned
  \item Whether the CBM scoring call happened after axis scoring (pipeline order verifiable)
\end{itemize}
\end{adminbox}

\begin{adminbox}{LLM Call Log — Prompt Verification}
The LLM Calls page shows every Ollama invocation with a 280-character prompt preview,
response preview, and token counts.
\begin{itemize}[leftmargin=1em]
  \item An auditor can verify that concept-scoring prompts list the expected concepts
        and request JSON output — confirming the CBM is using the LLM as a concept
        scorer, not as an end-to-end scorer.
  \item Prompt hashes (SHA-256) detect prompt drift: if a code change modifies a template,
        old vs.\ new calls are immediately distinguishable.
  \item Token count anomalies flag truncated context: a concept-scoring call with half the
        expected input tokens indicates the profile was silently trimmed.
\end{itemize}
\end{adminbox}

\begin{adminbox}{Daemon Activity and Evidence Tier Verification}
The Daemon Activity page logs every external observation: competitor pages checked,
grant calls found, legal entries deduplicated, KB staleness checks.
\begin{itemize}[leftmargin=1em]
  \item This verifies that T3 evidence tier upgrades ($m_i = 1.2$) are actually occurring
        and that daemon-observed data is flowing into profiles, raising multipliers in
        the scoring formula.
  \item Without this visibility, the evidence tier claim would be unverifiable to an
        auditor who cannot inspect Docker logs.
\end{itemize}
\end{adminbox}

\begin{adminbox}{Live Log Stream During a Diagnostic}
The Logs page tails the orchestrator ring-buffer at DEBUG level during a live run.
A judge can directly observe:
\begin{itemize}[leftmargin=1em]
  \item Wave ordering: Wave 0 completing before Wave 1 starts
  \item Per-axis score and breakdown as each axis service responds
  \item CBM bottleneck concept and projected lift as computed
  \item Anomaly flags as they fire
\end{itemize}
This live transparency is itself an evaluation demonstration: the system explains its
reasoning in real time to any observer with admin access.
\end{adminbox}

%% ─────────────────────────────────────────────────────────────────────────────
\section{Results Summary}
%% ─────────────────────────────────────────────────────────────────────────────

\begin{center}\small
\begin{tabular}{@{}llllll@{}}
\toprule
\textbf{Test} & \textbf{Tier} & \textbf{Metric} & \textbf{Target} & \textbf{Result} & \textbf{Status} \\
\midrule
Scoring determinism      & T2a & Identity across 10 runs              & 100\%       & 100\%       & \ppass \\
Rubric stability         & T2b & $\sigma \leq 0.15$ per field         & 0.15        & 0.07*       & \ppending \\
Anomaly recall           & T2c & 10/10 contradictions                 & 100\%       & 100\%       & \ppass \\
Maturity Macro-F1        & T1  & $\geq 0.65$ (5-class)               & 0.65        & 0.71*       & \ppending \\
Maturity Top-2 accuracy  & T1  & $\geq 0.85$                         & 0.85        & 0.92*       & \ppending \\
RAG Recall@3             & T3  & $\geq 0.80$                         & 0.80        & 0.90*       & \ppending \\
RAG MRR                  & T3  & $\geq 0.70$                         & 0.70        & 0.83*       & \ppending \\
Probe Recall gain        & T3  & vs.\ plain RRF                      & $+$         & $+$20 pp    & \ppass \\
Bottleneck stability     & CBM & Same concept, 10 runs               & $\geq$9/10  & 9/10        & \ppass \\
Weight convergence       & CBM & $\|\Delta\hat{\mathbf{w}}\| < 0.01$ & $<0.01$     & $<0.01$     & \ppass \\
\bottomrule
\end{tabular}
\end{center}

{\small $^*$\,Preliminary result on partial dataset.
Full T1 (50 profiles) and T3 (30 queries) datasets in preparation.}

\medskip
All T2 property tests and CBM stability checks pass.
Every preliminary T1 and T3 result exceeds its target on the partial dataset.
The \textbf{probe's $+20$ pp Recall improvement} is the most concrete quantitative
validation: it directly measures the axis direction probe's contribution against
the plain RRF baseline on real Tunisian KB queries.

\end{document}
